% !TEX root = ../main.tex
\chapter{Newtonian spacetime is curved!}\label{ch:c09}
\begin{paracol}{2}

\begin{sources}
\begin{tabularx}{\linewidth}{@{}l l L@{}}
\toprule
\textbf{Source} & \textbf{Length} & \textbf{Content}\\
\midrule
\seg{9}{1}{Lecture 9: Newtonian spacetime is curved!} & 1:48:32 & Newton's first two laws, the first as a measurement prescription for straight lines; why gravity cannot stay a force; the weak equivalence principle (the evacuated tube); Laplace's question and why curving space does not work; the full wisdom of Newton 1 (uniform motion) and world lines; the chart computation that turns Newtonian gravity into $\Gamma^\alpha{}_{00}=-F^\alpha/m$, the curvature (tidal tensor), Poisson's equation as the first Einstein equation, the acceleration as a sum; the geometric definition of a Newtonian spacetime (absolute time function, absolute space, uniform flow of time), future-directed vectors, the two laws in the new language; stratified atlases; Newton 2 in a chart, Coriolis and centrifugal terms. Student questions throughout, on torsion, on the rocket and on the Lagrangian.\\
\bottomrule
\end{tabularx}

\smallskip
\textbf{Transcript} \file{materials/transcripts/C09-V01.txt} (YouTube captions). There is no tutorial for this lecture: the tutorial of lectures 9 and 10 is in chapter~\ref{ch:c10}.
\end{sources}

\begin{reading}[Where this chapter sits in the other notes]
\emph{Theoretische Mechanik} notes (Schuller's German classical-mechanics lectures, which the lecture advertises at the end): chapters 8 and 9 there treat Newtonian kinematics and absolute space [titles not checked], and chapter 6 the connections used here (chapter~\ref{ch:c07} here). The auto-parallel equation, the connection coefficient functions and the transformation law used all over this chapter are from lecture 7 (chapter~\ref{ch:c07}); the curvature and the Ricci tensor are from lecture 8 (chapter~\ref{ch:c08}, not yet written), and are only used here as results.
\end{reading}

% ==================================================================
\section{Newton's laws and the problem with gravity}\label{sec:c09-newton}

Welcome to the ninth lecture \ts{9}{1}{0:00}. The title, \emph{Newtonian spacetime is curved}, ``may come as a shock'', because one probably thought it is not curved, and in the standard formulation it indeed is not. The question is whether this standard formulation holds up when the axioms are inspected properly: the point is that Newtonian spacetime must be considered curved, and that this curvature \emph{absorbs the effect of the gravitational force}, so that gravity can no longer be considered a force. Today is a physics lecture, but much of what was learned (especially the auto-parallel equation) is used
\ts{9}{1}{0:21}\ts{9}{1}{0:42}\ts{9}{1}{1:02}.

\begin{definition}[Newton's first two laws, as worded in the lecture]\label{def:c09-newton}
Newton wrote them with little drawings and words, ``very carefully worded'':
\begin{enumerate}
  \item (Newton 1) A body on which no force acts moves uniformly along a straight line.
  \item (Newton 2) The deviation of a body's motion from such uniform straight motion (today: the acceleration) is effected by a force, and it is reduced by a factor, the reciprocal of the body's \emph{inertial mass}: the bigger the inertial mass, the less the same force accelerates the body.
\end{enumerate}
\end{definition}
\ts{9}{1}{1:23}\ts{9}{1}{1:45}\ts{9}{1}{2:05}\ts{9}{1}{2:25}.
Usually one adds ``in an inertial system'' to both; the lecture does not, because an inertial system has to do with choosing coordinates, whereas the statement is \emph{absolute}: one knows what a straight line is independently of the reference system (or at least one can formulate this today)
\ts{9}{1}{1:45}.

\paragraph{The first law as a measurement prescription.} Read as a postulate of what a particle does under various circumstances, the first law is merely the special case $F=0$ of the second (``Newton was not that naive, he did not need to emphasise the special case in an extra axiom''). To have an important role it must be read differently: one assumes that, for some reason, one \emph{knows} that no force acts, and uses such particles to test experimentally what a straight line is. Imagine outer space with absolutely no force: one throws a stone, and the first law instructs one to say ``that is a straight line'': one is testing the geometry of space.
So \textbf{remark 1: the first axiom, to be relevant, must be read as a measurement prescription for the geometry of space} (``and actually we will add something behind space later''). This makes sense: in science one had better define an entity like ``straight line''; on paper in mathematics it is obvious what a straight line is, but is it obvious what a straight line is in nature? Newton recognises this and makes of it a measurement prescription. The second law then tells what happens to the other particles on which a force acts, for instance because they have charges and react to some fields, ``a modern Faraday point of view''
\ts{9}{1}{2:47}\ts{9}{1}{3:07}\ts{9}{1}{3:27}\ts{9}{1}{3:48}\ts{9}{1}{4:08}\ts{9}{1}{4:30}\ts{9}{1}{4:52}\ts{9}{1}{5:20}\ts{9}{1}{5:41}\ts{9}{1}{6:01}.

\paragraph{The problem.} There is a second important fact observed in nature: \emph{gravity affects all massive particles}. With a single massive particle in the universe (neglecting its gravitational self-interaction) the first law might be realisable; but as soon as there is a second particle, and one thinks of gravity as a force, which is the standard view, one can no longer use the first particle to measure the geometry of space: in Newtonian theory all particles have mass and all masses gravitate at arbitrarily large distances, so the assumption ``no force acts'' can no longer be applied. ``We say, come on, we never use this anyway and it works fine'': but the axiom is there, and it is very important to say what a straight line is, because \emph{if one does not know what a straight line is one does not know what the deviation from a straight line is}.
Being educated, one knows that a straight line is determined by a covariant derivative, i.e.\ by the $\Gamma$'s: with $V$ the tangent vector of a curve, a straight line is $\nabla_VV=0$ (the auto-parallel equation), and the deviation from it, $\nabla_VV\neq0$, is the acceleration, a vector again. Without knowing the straight lines one does not know the $\Gamma$'s and cannot calculate the acceleration (``we will do this more precisely'')
\ts{9}{1}{6:23}\ts{9}{1}{6:43}\ts{9}{1}{7:03}\ts{9}{1}{7:24}\ts{9}{1}{7:44}\ts{9}{1}{8:05}\ts{9}{1}{8:25}.
Conclusion: since gravity universally acts on every particle, in a universe with at least two particles \emph{gravity must not be considered a force} if Newton 1 is supposed to remain applicable. This is a serious problem, because one does not know how to encode gravity other than as a force \ts{9}{1}{8:46}\ts{9}{1}{9:10}\ts{9}{1}{9:34}\ts{9}{1}{10:03}.

\section{Laplace's question}\label{sec:c09-laplace}

The first to think about this was Laplace (1749 to 1827, as said in class): ``can gravity be encoded in a curvature of space, such that its effects show if particles, under the influence of no other force, are postulated to move along straight lines in this curved space?'' (One still thinks of gravity as a force, but wants to get rid of it.)
\begin{history}[Laplace's question]
The lecture dates Laplace as 1749 to 1827 and says that he spent quite a while (``several days'') thinking about this; neither the dates nor this are checked here. (As said in class, not checked.)
\end{history}
In some science museums one sees a bowl-shaped throat on which one puts a coin and flips it: it runs around and spirals inwards because of friction; without friction it would go around on an orbit, and with less energy on a smaller one. ``Could this be the motion of the planets around the Sun, one dimension higher? That is the rough idea''
\ts{9}{1}{10:24}\ts{9}{1}{10:45}\ts{9}{1}{11:08}\ts{9}{1}{11:31}\ts{9}{1}{11:53}\ts{9}{1}{12:16}\ts{9}{1}{12:37}\ts{9}{1}{12:57}.
The answer to Laplace's question is a ``resounding no'': it is absolutely impossible to come up with a curved space such that this works, and it is a two-liner. One does it nevertheless because it is instructive to see what does not work, as it can be repaired \ts{9}{1}{13:18}\ts{9}{1}{13:38}.

\paragraph{The standard view: gravity as a force.} Newton 2 with the gravitational force field: mass times acceleration equals the force, $m\,\ddot x=F(x)$ (at a point, the Euler picture; ``this is not the acceleration, but let us not talk about it right now, that comes later''). The gravitational force field satisfies Poisson's equation: writing $F$ as minus the gradient of the potential $\Phi$ times mass, $\Delta\Phi=4\pi G\rho$ with $\rho$ the mass density of matter, so Poisson's equation is the gravitational field equation: depending on the masses around, it determines the force law, which one plugs into Newton 2 to see how a particle reacts.
Now the thing: the gravitational force is proportional to the mass of the body that produces it, \emph{and} to the mass of the body on which it acts (``very special''), so one can cancel the mass of the body on which it acts altogether: $\ddot x=f(x)$ with $F=mf$, and no more mention of the mass of the accelerated body \ts{9}{1}{13:38}\ts{9}{1}{14:00}\ts{9}{1}{14:21}\ts{9}{1}{14:41}\ts{9}{1}{15:02}\ts{9}{1}{15:22}\ts{9}{1}{15:42}\ts{9}{1}{16:03}.

\paragraph{The weak equivalence principle, and the experiment.} Is this true? ``Theory does not tell us, let us go to experiment''; in theory it is called the \textbf{weak equivalence principle} (the equality of the inertial mass of Newton 2 and the gravitational mass), ``but we can give it names all day long, experiment decides''. In a tube there are a little feather and a little bead; erected vertically, everyday experience says that the bead falls like that and the feather falls more slowly (``less accelerated, maybe because you think it has less mass''); but this can all be attributed to the air resistance in the tube. Then the tube is evacuated (live in the lecture hall) and the experiment repeated: it is not true that the bead is accelerated more strongly than the feather, they are both accelerated the same, within the error of measurement. (``Tada.'') One gives this the grand name of weak equivalence principle, ``which is then great for philosophers without experiments''
\ts{9}{1}{16:24}\ts{9}{1}{16:51}\ts{9}{1}{17:11}\ts{9}{1}{17:31}\ts{9}{1}{17:54}\ts{9}{1}{18:15}\ts{9}{1}{19:08}\ts{9}{1}{19:28}\ts{9}{1}{22:14}.

\begin{coursenote}[Student questions during and after the demonstration]
\begin{itemize}
  \item \emph{Why not just dispose of Newton 1?} A student sees two ways out of the problem: keep Newton's first law, or drop it; but then ``what do you mean by deviation from a straight line'', and one could refer to the straight lines of Euclidean geometry; then why does physical space feature a Euclidean geometry, and at that point one reinstates the law and sees whether it can be saved: ``it is a fully conservative approach'', nothing said here goes against anything written in Newton's \emph{Principia}, only against what was made of it afterwards by using vectors \ts{9}{1}{19:49}\ts{9}{1}{20:10}\ts{9}{1}{20:32}\ts{9}{1}{20:52}.
  \item \emph{Angular momentum on a curved bowl.} A second question, the example usually brought up for the conservation of angular momentum: if one sits on the particle going around, why can one not say that it is a free motion (``if you attach an additional frame and you return, the frame still points in the same direction'')? This is the free-falling motion; one might think that one can feel that one goes around, like a centrifugal force or so, ``we will come to this'' (see the last sections) \ts{9}{1}{20:52}\ts{9}{1}{21:13}\ts{9}{1}{21:33}\ts{9}{1}{21:54}.
\end{itemize}
\end{coursenote}

\paragraph{Laplace's idea as an auto-parallel equation.} After cancelling $m$, Laplace wants to interpret the equation $\ddot x=f$ not as acceleration and force, but as the equation of a straight line: ``Laplace does not say: change the equation, put an extra term; he says: let us interpret this equation as an auto-parallel equation, because that is the equation for a straight line''. Taking $f$ to the other side, the question becomes whether it can be written for $x^\alpha(t)$, $\alpha=1,2,3$ (three-dimensional space), as
\[
  \ddot x^\alpha(t)+\Gamma^\alpha{}_{\beta\gamma}\bigl(x(t)\bigr)\,\dot x^\beta(t)\,\dot x^\gamma(t)=0,
\]
that is, as an auto-parallel with respect to some $\nabla$ on space. The answer is \emph{no}: $\Gamma$ can only depend on the position $x$, but the equation has two velocity factors, which cannot be matched (not for zero velocity, and not for any velocity: with double the velocity one gets four times the term; whatever choice one makes for the $\Gamma$'s in terms of $f$, one cannot get rid of the effect of the $\dot x$'s; even if $\Gamma$ depended on $\dot x$ it would be very difficult to get rid of them). So \textbf{one cannot find $\Gamma$'s such that Newton's equation for the gravitational force takes the form of an auto-parallel equation}: the equation does not have that form. The standard example from the museums is ``entirely wrong, you cannot curve space that way''. (Laplace found it without the auto-parallel equation: ``he did not have it, he spent several days thinking and realised there are other inconsistencies'')
\ts{9}{1}{22:35}\ts{9}{1}{22:55}\ts{9}{1}{23:16}\ts{9}{1}{23:37}\ts{9}{1}{23:58}\ts{9}{1}{24:18}\ts{9}{1}{24:39}\ts{9}{1}{25:02}\ts{9}{1}{25:23}\ts{9}{1}{25:43}.

\begin{coursenote}[Student question: can the $\Gamma$ be made to vanish along the fall?]
A student recalls (from lecture 7) that the $\Gamma$'s can be made to vanish at a point, and asks whether this could be done along the whole curve. Schuller: at one point yes, by going to the coordinate system that falls with the chalk; to do it all along the curve one must ``climb up here and jump down'' and fall with it. The problem is the Earth being round: someone at the North Pole making the experiment can transform away the gravitational force where he is in his falling coordinate system, but at the same moment someone at the South Pole in the system of the falling guy up here will not have the effect transformed away. A \emph{homogeneous} gravitational field can be transformed away entirely (more later), an inhomogeneous one cannot
\ts{9}{1}{26:06}\ts{9}{1}{26:26}\ts{9}{1}{26:47}\ts{9}{1}{27:09}\ts{9}{1}{27:29}\ts{9}{1}{27:50}.
\end{coursenote}

\section{The full wisdom of Newton 1: world lines}\label{sec:c09-worldlines}

As so often, Laplace did not look carefully. What Laplace asked was a curvature of space such that particles move along straight lines; he did not read Newton 1 properly: it talks not only about a straight line but about \emph{uniform motion along a straight line}. ``The full wisdom of Newton 1 is to use also the information that particles under the influence of no force move uniformly.'' Laplace's idea implemented only the first half; ``as is often said in research proposals, it is a beginning'', but one has to take everything into account \ts{9}{1}{27:50}\ts{9}{1}{28:11}\ts{9}{1}{28:33}\ts{9}{1}{28:54}\ts{9}{1}{29:35}\ts{9}{1}{29:56}.

A drawing in a two-dimensional flat space: uniform motion with equal amounts of time, a straight line; a second straight line traversed non-uniformly, getting slower and slower. Something very difficult to explain to students at first: \emph{a curve is more than the set of its points: it is the set of its points and its parametrisation}. The way to capture the difference is a data structure everyone knows from school: one draws a space-time diagram (for a harmonic oscillator, how does the position depend on time?). In a one-dimensional world a straight uniform motion becomes a straight line in the diagram; a straight non-uniform motion (in one dimension motion is always straight; the point is the difference between straight and uniform) becomes a curve. ``We attach a time axis and we do not introduce new data, we just represent the data in a different fashion''
\ts{9}{1}{30:16}\ts{9}{1}{30:37}\ts{9}{1}{31:01}\ts{9}{1}{31:21}\ts{9}{1}{31:43}\ts{9}{1}{32:04}\ts{9}{1}{32:24}.
\begin{coursenote}[A slip at the board]
Drawing equal time intervals Schuller gets confused for a moment about whether the non-uniform motion is getting slower or faster in the diagram (``as a relativist this is a rather embarrassing moment, we are trying to cut it off the video''; there is no excuse of ordering the axes differently in relativity, ``because there are no axes''); the conclusion is that in equal times it makes less distance, i.e.\ slower \ts{9}{1}{32:45}\ts{9}{1}{33:05}\ts{9}{1}{33:25}.
\end{coursenote}

The insight: in spacetime, if one attaches time to space and considers it an entity, one can \emph{ignore the parametrisation} of the curve (``nobody cares whether you parametrise this curve in spacetime in the most crazy manner, because for every point where you are you just look at the time axis and you know what time it is''): the information has been stored in the extra coordinate. And better: ``uniform and straight motion in space is simply \emph{straight} motion in spacetime'', no further questions asked. So Laplace's idea might work if one only looks at straight lines, but \emph{in Newtonian spacetime}; ``this has nothing to do with relativity, it is just drawing what you already drew in school for the harmonic oscillator, just representing the data in a different fashion, one to one''. Then one can use the ideas about straight motion, otherwise one would have to invent new differential geometry implementing straight uniform motion (``it can be done too, but not in that way'')
\ts{9}{1}{33:46}\ts{9}{1}{34:07}\ts{9}{1}{34:28}\ts{9}{1}{34:51}\ts{9}{1}{35:12}\ts{9}{1}{35:33}\ts{9}{1}{35:53}.

\paragraph{World lines.} Let $x\colon\R\to\R^3$ be a particle trajectory in space. One converts it back and forth into a \textbf{world line} (the history of the particle: from it one knows its entire story) by assigning to each $t$ the quadruple of real numbers
\[
  X(t):=\bigl(t,\,x^1(t),\,x^2(t),\,x^3(t)\bigr)
\]
reproducing what one already had. Picture: the Sun (a very massive body) and a small body (``me'') on a bound orbit in space, with Kepler's law of equal areas in equal times, so that the parametrisation plays a role; in spacetime the circular or elliptical motion becomes a \emph{spiral}, and from the spiral one gets the orbit back by projecting, forgetting the first component. ``There is nothing happening, no new physics, it is just a different way to represent the history of a particle: that is all it takes to make Laplace's idea work, but in spacetime'' (a joke about two popes meeting is told in spacetime but ``never works'') \ts{9}{1}{36:16}\ts{9}{1}{36:38}\ts{9}{1}{36:58}\ts{9}{1}{37:21}\ts{9}{1}{37:42}\ts{9}{1}{38:04}\ts{9}{1}{38:25}\ts{9}{1}{38:46}\ts{9}{1}{39:06}\ts{9}{1}{39:26}\ts{9}{1}{39:47}.

\section{The chart computation: gravity is $\Gamma^\alpha{}_{00}$}\label{sec:c09-chart}

Let the small $x$ satisfy Newton's gravitational law, with the mass already cancelled: $\ddot x^\alpha=f^\alpha(x)$, $\alpha=1,2,3$ (Greek indices run from $1$ to $3$, Latin indices from $0$ to $3$). The following are ``trivial rewritings'', all in charts. Let $X^0$ denote the zeroth component (a manner of numbering): $X^0(t)=t$, hence
\[
  \dot X^0=1,\qquad\ddot X^0=0,
\]
and for the spatial components, since $X^\alpha(t)=x^\alpha(t)$,
\[
  \ddot X^\alpha-f^\alpha(X)=0
\]
(strictly a different function, $f^\alpha$ regarded as a function on spacetime which ignores its first entry; ``that is not very deep''). The first equation is simply true and does not cause trouble. Now multiply the terms of the second by $\dot X^0\dot X^0=1\cdot1$, which changes nothing. Then the two equations together are \emph{perfectly equivalent} to the set, $a=0,1,2,3$,
\[
  \ddot X^a+\Gamma^a{}_{bc}\,\dot X^b\dot X^c=0
\]
for suitable $\Gamma$'s: an auto-parallel equation \emph{in Newtonian spacetime} \ts{9}{1}{40:08}\ts{9}{1}{40:33}\ts{9}{1}{40:53}\ts{9}{1}{41:15}\ts{9}{1}{41:35}\ts{9}{1}{41:56}\ts{9}{1}{42:17}\ts{9}{1}{42:38}\ts{9}{1}{42:58}\ts{9}{1}{43:20}\ts{9}{1}{43:41}.

\begin{proposition}[Newtonian gravity is an auto-parallel equation in spacetime]\label{prop:c09-gamma}
In the chart $(t,x^1,x^2,x^3)$ choose the connection coefficient functions
\[
  \Gamma^\alpha{}_{00}=-f^\alpha,\qquad\text{all other }\Gamma^a{}_{bc}=0.
\]
Then the auto-parallel equation $\ddot X^a+\Gamma^a{}_{bc}\dot X^b\dot X^c=0$ with $X^0=t$ is Newton's equation $\ddot x^\alpha=f^\alpha(x)$.
\end{proposition}
\begin{proof}
Compare term by term. For $a=0$: $\ddot X^0=0$ has no $\Gamma$-term, so $\Gamma^0{}_{bc}=0$ for all $b,c$ (otherwise there would be terms that Newton's equation does not have). For $a=\alpha$: if any $\Gamma^\alpha{}_{\beta\gamma}$ with $\beta,\gamma$ spatial were non-zero, it would produce a term $\dot X^\beta\dot X^\gamma$ (say $\dot X^1\dot X^1$), which does not occur; if one of the two lower indices is $0$ and the other spatial one gets a term $\dot X^0\dot X^\beta=\dot X^\beta$ that does not occur either. So all $\Gamma^\alpha{}_{\beta\gamma}$ and $\Gamma^\alpha{}_{0\beta}$, $\Gamma^\alpha{}_{\beta0}$ vanish. What remains is $\Gamma^\alpha{}_{00}\dot X^0\dot X^0=\Gamma^\alpha{}_{00}$, which must equal $-f^\alpha$. ``Now it makes sense with the indices''
\ts{9}{1}{44:04}\ts{9}{1}{44:25}\ts{9}{1}{44:46}\ts{9}{1}{45:07}\ts{9}{1}{45:28}\ts{9}{1}{45:50}.
\end{proof}
So in the chart one started with (``the Cartesian chart, whatever it is called: we play the game for a time being until we made our point'') one can choose the $\Gamma$'s to be minus the gravitational force per mass and all others zero \ts{9}{1}{45:50}\ts{9}{1}{46:10}.

\paragraph{Is it a coordinate artifact?} One may ask whether the effect could be transformed away, since this is just one chart: one could use the transformation law of the $\Gamma$'s, and the $\Gamma$'s can be non-zero in some charts and zero in others, depending on whether the \emph{curvature} vanishes or not (lecture 8). The answer is no, since one can calculate the Riemann curvature of the connection defined this way (chapter~\ref{ch:c08}); in this chart the \emph{only} non-vanishing components are
\[
  R^\alpha{}_{0\beta0}=-\frac{\partial f^\alpha}{\partial x^\beta},
\]
the \textbf{tidal force tensor} (minus the Hessian of the gravitational potential, up to sign conventions): ``if you study how tidal forces act, you look at this guy''. Contracting the first and the third index (``it was on the last tutorial sheet'') gives the Ricci tensor, whose only non-vanishing component is
\[
  R_{00}=-\frac{\partial f^\alpha}{\partial x^\alpha}=4\pi G\rho,
\]
by Poisson's equation ($-\partial_\alpha f^\alpha=4\pi G\rho$, since $f=-\nabla\Phi$)
\ts{9}{1}{46:31}\ts{9}{1}{46:51}\ts{9}{1}{47:15}\ts{9}{1}{47:35}\ts{9}{1}{47:56}\ts{9}{1}{48:17}\ts{9}{1}{48:38}\ts{9}{1}{48:59}\ts{9}{1}{49:19}.
\begin{coursenote}[Conventions]
The values $R^\alpha{}_{0\beta0}=-\partial_\beta f^\alpha$ and $R_{00}=-\partial_\alpha f^\alpha$ follow from $R^a{}_{bcd}=\partial_c\Gamma^a{}_{db}-\partial_d\Gamma^a{}_{cb}+\Gamma\Gamma-\Gamma\Gamma$ and $R_{bd}=R^a{}_{bad}$ with the $\Gamma$'s of \cref{prop:c09-gamma}: only $\partial_\beta\Gamma^\alpha{}_{00}$ survives, since all the products of $\Gamma$'s vanish. (Not said in class; the curvature convention is the one of the next lecture [not yet checked, chapter~\ref{ch:c08} is not written].)
\end{coursenote}

\paragraph{The first Einstein equation.} Interesting if one knows general relativity: one could call $\rho$ the $00$ component of an energy-momentum tensor $T_{ab}$, and one gets something like $R_{00}=8\pi G\,T_{00}$ (``that is now really in a sense cheating'': the factors $2$ and $4\pi$ versus $8\pi$ are fudged on the board: $R_{00}=4\pi G\rho$, ``no, $8\pi G$''), the first of the ten gravitational equations postulated by Einstein in 1912. In relativistic physics Einstein said $R_{ab}=8\pi G\,T_{ab}$; ``this is not entirely correct but almost correct: spacetime Ricci curvature is the energy-momentum of matter''. Poisson's law says that for the curvature of spacetime in Newtonian gravity $R_{00}$ is proportional to $\rho$; for dust (many non-interacting particles with only mass) this looks exactly like it. So Poisson's equation, which everyone knows, makes one understand the basic idea of Einstein's equations, at least in the old, not-correct form (an extra term is missing that makes it correct, ``and that has to do with relativity''): the field equation of gravity is right here
\ts{9}{1}{49:40}\ts{9}{1}{50:01}\ts{9}{1}{50:23}\ts{9}{1}{50:48}\ts{9}{1}{51:08}\ts{9}{1}{51:29}\ts{9}{1}{51:50}.
\begin{coursenote}[Factors in the Einstein equation]
(Not said in class.) The correct field equation is $R_{ab}=8\pi G\bigl(T_{ab}-\tfrac12Tg_{ab}\bigr)$, and for dust, $T_{ab}=\rho u_au_b$, in a signature with $g_{00}=-1$, this gives $R_{00}=8\pi G(\rho-\tfrac12\rho)=4\pi G\rho$: exactly Poisson's equation. This is the ``extra term'' that the lecture mentions without writing it, and the reason for the board's factor of $2$. Einstein's 1912 date is as said in class, not checked.
\end{coursenote}

Because $R_{00}$ is non-zero in the presence of masses, the curvature is non-zero, and by the theorem of lecture 7/8 \emph{there is no coordinate system in which all the $\Gamma$'s vanish}: with the Earth, an inhomogeneous gravitational field, there is no system in which for every particle one can make them vanish. ``It is not a coordinate artifact, it is not `oh, choose better coordinates and the effect is gone': wouldn't it be funny if it was a coordinate artifact, because coordinates are in our mind but you feel whether you accelerate or not''. \textbf{Conclusion: Laplace's idea works on Newtonian spacetime}
\ts{9}{1}{51:50}\ts{9}{1}{52:11}\ts{9}{1}{52:32}\ts{9}{1}{52:53}\ts{9}{1}{53:15}.

\paragraph{It is spacetime curvature, not space curvature.} The $\Gamma$'s that are non-zero have time indices: all components $R^\alpha{}_{\beta\gamma\delta}$ with spatial indices vanish, so \emph{space} is not curved: ``spacetime is not curved in the space directions, roughly speaking; it is curved in the extra direction together with space''. Only $R^\alpha{}_{0\beta0}$ does not vanish, so not all components of the Riemann tensor vanish \ts{9}{1}{53:15}\ts{9}{1}{53:35}\ts{9}{1}{53:57}\ts{9}{1}{54:17}.

\paragraph{The acceleration is a sum.} A student points out that the right-hand side of the equation certainly transforms as a tensor, but $\Gamma$ does not (nor does $\ddot X$): what about the transformation behaviour of the left-hand side $\ddot X^a+\Gamma^a{}_{bc}\dot X^b\dot X^c$? ``Very good question.'' The velocity components transform as a vector. If one takes the covariant derivative of the velocity of a curve in the direction of the velocity, its $a$-th component is exactly this expression; it is a vector field along the curve, the \textbf{acceleration} $a=\nabla_VV$ (written with a fancy $a$ because of the index $a$). \emph{The acceleration is a vector}, but the second derivative of the components of a curve is not the components of that acceleration: $\ddot X^a$ does not transform as vector components, $\Gamma\dot X\dot X$ does not either, ``in the opposite way'' to the inhomogeneous term in the transformation of the $\Gamma$'s, and \emph{only the sum} of the two cancels the troublesome terms. One cannot say ``that is the acceleration and that is an extra acceleration'': together they are the acceleration, and of course it is chart independent (``if you sit on a rocket you feel whether you accelerate or not; how crazy would it be if somebody thinks another coordinate system and you feel something different?'')
\ts{9}{1}{54:17}\ts{9}{1}{54:38}\ts{9}{1}{55:02}\ts{9}{1}{55:23}\ts{9}{1}{55:44}\ts{9}{1}{56:05}\ts{9}{1}{56:26}\ts{9}{1}{56:46}\ts{9}{1}{57:07}\ts{9}{1}{57:27}.

\paragraph{Newton 1 in spacetime.} So Newtonian theory is reformulated such that gravity is a curvature of Newtonian spacetime, and the first axiom takes the form: \emph{a particle under the influence of no other force (gravity no longer counts as one) has a world line that is straight}, ``end of message'': not ``uniformly straight in space'', but \emph{straight in spacetime}. If two particles are dropped, initially parallel lines start to converge (``small Earth''): this is another way to identify curvature, \textbf{geodesic deviation}, as on the sphere where lines going parallel from the equator finally meet. So the problem is solved: space is nice and flat, spacetime is curved, and the first axiom can be put into practice. ``Case closed'', up to the terribly chart-dependent formulation, all with components \ts{9}{1}{57:27}\ts{9}{1}{57:49}\ts{9}{1}{58:09}\ts{9}{1}{58:30}\ts{9}{1}{58:51}\ts{9}{1}{59:12}.

\section{The geometric formulation: Newtonian spacetime}\label{sec:c09-geometric}

``Adequate sketch of the foundations of the geometric formulation of Newton's axioms'': no new physics, ``just what Newton wrote down in words, implemented properly with the mathematics that makes it work, rather than the mathematics that seems to make it work''. A fresh start (the lecture says: forget everything from before; it was only to give an idea how one gets here) \ts{9}{1}{59:12}\ts{9}{1}{59:38}\ts{9}{1}{59:58}.

\begin{definition}[New concepts: Newtonian spacetime, absolute time function]\label{def:c09-nst}
A \textbf{Newtonian spacetime} is a quintuple $(M,\mathcal O,\mathcal A,\nabla,t)$, where
\begin{itemize}
  \item $(M,\mathcal O,\mathcal A)$ is a four-dimensional smooth manifold (``it has no direction: this is space and this is time''; the danger of the word is to misspell it as ``space plus time''; spacetime is one thing in which at first one cannot distinguish between space and time, like the sphere, a two-dimensional thing in which one does not say where is one direction over the other; one gets a derived notion of space and time only under certain circumstances: ``delete all your memory'');
  \item $\nabla$ is a covariant derivative on $M$ (``Newton talks about straight lines, uniform straight lines; we translate this into straight lines in spacetime'');
  \item $t\colon M\to\R$ is a smooth function, the \textbf{absolute time} (``Newton talks insistently about absolute time''),
\end{itemize}
satisfying the following axioms, which translate what Newton says:
\begin{enumerate}
  \item (\emph{absolute space}) $\mathrm dt\neq0$ everywhere: the gradient of the absolute time function does not vanish anywhere on $M$;
  \item (\emph{absolute time flows uniformly}) $\nabla\mathrm dt=0$: for every vector field $X$, $\nabla_X\mathrm dt$ is the zero covector field;
  \item (added at the end of the lecture, see below) $\nabla$ is torsion free.
\end{enumerate}
\end{definition}
\ts{9}{1}{1:00:19}\ts{9}{1}{1:00:39}\ts{9}{1}{1:01:00}\ts{9}{1}{1:01:20}\ts{9}{1}{1:01:41}\ts{9}{1}{1:02:01}\ts{9}{1}{1:02:22}\ts{9}{1}{1:02:44}\ts{9}{1}{1:07:00}\ts{9}{1}{1:24:44}\ts{9}{1}{1:25:04}.

\paragraph{Absolute space.} Newton says there is absolute space; the structure provides only an absolute time, but absolute space can be derived. The requirement is that the gradient of $t$ is non-zero everywhere. Define \emph{absolute space at time $t_0$} as
\[
  S_{t_0}:=\{p\in M\mid t(p)=t_0\},
\]
the collection of all points (events) with the time label $t_0$ (every event is given a universal clock time). Because $\mathrm dt\neq0$ the manifold is the \emph{disjoint union} of these spaces, it is \textbf{foliated} into $S_{t_1},S_{t_2},\dots$ (three-dimensional picture of a four-dimensional spacetime); the universal clock gives the same time for all points of one such space. If the gradient vanished somewhere, level lines of the time function would intersect or merge there: a problem, so it is excluded. ``So this innocent requirement on the absolute time yields a notion of absolute space, and of course this is how we think of spacetime: spaces stacked together''; one might ask why they are not stacked in a wobbly way (``have we talked about shape yet? No, we have not'')
\ts{9}{1}{1:02:44}\ts{9}{1}{1:03:07}\ts{9}{1}{1:03:28}\ts{9}{1}{1:03:49}\ts{9}{1}{1:04:11}\ts{9}{1}{1:04:32}\ts{9}{1}{1:04:52}\ts{9}{1}{1:05:13}\ts{9}{1}{1:05:35}\ts{9}{1}{1:05:57}\ts{9}{1}{1:06:18}\ts{9}{1}{1:06:38}.
\begin{coursenote}[Student question: what if the gradient vanishes?]
A student asks whether a point where the gradient vanishes would correspond to time stopping. Yes, ``time would stop here for a while, that is what the gradient is; and that is of course meaningless, not what we want in physics, but it could be if we just say it is a smooth function'': the axiom excludes it \ts{9}{1}{1:08:48}\ts{9}{1}{1:09:08}.
\end{coursenote}

\paragraph{Time flows uniformly.} ``Newton says: absolute time flows uniformly''; what does it mean? ``One second per second, what is that supposed to mean?'' If one wants to be clear: the covariant derivative of the gradient vanishes everywhere, $\nabla_X\mathrm dt=0$ for any $X$, written as a field equation. In a stratified atlas (below) one sees that it means that time flows uniformly. The two axioms are a compatibility axiom since they relate two things Newton talks about ``so eloquently''; what exactly they mean is sometimes a matter of heated debate, ``this is a clear formulation of what it means, and you can make this work''
\ts{9}{1}{1:06:38}\ts{9}{1}{1:07:00}\ts{9}{1}{1:07:20}\ts{9}{1}{1:07:42}\ts{9}{1}{1:09:29}.
\begin{coursenote}[Student question: the slot of $\nabla\mathrm dt$]
A student asks about the open slot in $\nabla_X\mathrm dt$. One could put $\partial/\partial x^a$ there in any chart and take the $b$-th component; because $\nabla\mathrm dt$ is $C^\infty$-linear in the empty slot, it is a $(0,2)$ tensor field if that slot is left unfilled, and the requirement is that it is the zero tensor field of that type: ``I maintain that this equation is good, an equal sign on the space of $(0,2)$ tensor fields'' (in coordinates a little later) \ts{9}{1}{1:07:42}\ts{9}{1}{1:08:02}\ts{9}{1}{1:08:23}.
\end{coursenote}

Schuller says that his main point for these lectures was made before: \emph{one can view gravity as a curvature of spacetime already in Newtonian mechanics}, not as an alternative formulation but because, reading the first axiom, ``what other choice do I have?''; this is the philosophical point of view of general relativity later, and it has all to do with gravity and nothing to do with relativity: ``it is not relativity that forces us to spacetime as far as gravity is concerned, it is gravity itself and the understanding of the relation between geometry and gravity'', the most important lesson for a theory builder
\ts{9}{1}{1:09:49}\ts{9}{1}{1:10:10}\ts{9}{1}{1:10:30}\ts{9}{1}{1:10:51}.

\begin{definition}[New concepts: future-directed, spatial and past-directed vectors]\label{def:c09-fut}
A vector $X\in T_pM$ in a Newtonian spacetime is called \textbf{future directed} if $\mathrm dt(X)>0$, \textbf{spatial} if $\mathrm dt(X)=0$, and \textbf{past directed} if $\mathrm dt(X)<0$. A \textbf{future-directed curve} is a curve whose velocity vectors are future directed at every point.
\end{definition}
\ts{9}{1}{1:11:12}\ts{9}{1}{1:11:33}\ts{9}{1}{1:11:55}.
Picture: the level surfaces of the absolute time function are the absolute spaces at the various times; a spacetime vector at a point of $S_1$ pointing to $S_2$ (time increases from $S_1$ to $S_2$) is future directed, because the gradient (a covector field) applied to it gives a positive number; a vector tangent to the level surface is spatial, since the gradient on a tangent vector of a level surface vanishes (``this old theorem''); a vector pointing the other way is past directed \ts{9}{1}{1:12:16}\ts{9}{1}{1:12:37}\ts{9}{1}{1:12:58}.

\begin{definition}[Newton 1 and Newton 2 in a Newtonian spacetime]\label{def:c09-newtongeom}
\begin{enumerate}
  \item (Newton 1) The world line of a particle under the influence of no force (gravity is not one in this view) is a \emph{future-directed auto-parallel} with respect to $\nabla$: with $V$ the velocity of the world line $X$,
        \[
          \nabla_VV=0,\qquad\mathrm dt(V)>0.
        \]
  \item (Newton 2) For a world line $X$ with velocity $V$,
        \[
          m\,\nabla_VV=F,
        \]
        where the force $F$ is a \emph{spatial} vector field, $\mathrm dt(F)=0$ (``a force is only supposed to accelerate you in the spatial directions''). The acceleration is $\nabla_VV$; ``nobody said $m\cdot a=F$ is wrong, I only said that the acceleration is not $\ddot X$''.
\end{enumerate}
\end{definition}
\ts{9}{1}{1:13:20}\ts{9}{1}{1:13:41}\ts{9}{1}{1:14:07}\ts{9}{1}{1:14:30}\ts{9}{1}{1:14:51}\ts{9}{1}{1:15:11}\ts{9}{1}{1:15:33}\ts{9}{1}{1:15:55}\ts{9}{1}{1:16:16}\ts{9}{1}{1:16:40}.
(Note the future-directedness: a particle runs into the future and not the past, in spacetime one must distinguish future and past.) ``No more funny talk''
\ts{9}{1}{1:14:51}.

\paragraph{Stratified atlases.} One more convention, ``incredibly practical'': restrict attention to \textbf{stratified atlases}, atlases whose charts $(U,x)$ have the property that the zeroth coordinate function coincides with the absolute time function, $x^0=t|_U$. One can require this: $t$ is defined on all of $M$ and can be restricted to a chart domain, while the other coordinates are only defined on $U$; ``if you are a geometer you can require on your charts whatever you like as long as they still cover the manifold; your charts are in your head''. All formulas are right in any atlas, they just become simpler in a stratified one, and nothing depends on it; in a non-stratified atlas one gets extra terms, but they are coordinate artifacts
\ts{9}{1}{1:17:00}\ts{9}{1}{1:17:21}\ts{9}{1}{1:17:43}\ts{9}{1}{1:18:04}\ts{9}{1}{1:18:24}\ts{9}{1}{1:18:45}\ts{9}{1}{1:19:06}.

In a stratified atlas, $\mathrm dt=\mathrm dx^0$ and the axiom $\nabla\mathrm dt=0$ reads, for all $a$,
\[
  0=\nabla_{\partial/\partial x^a}\,\mathrm dx^0=-\Gamma^0{}_{ba}\,\mathrm dx^b\quad\Longrightarrow\quad\Gamma^0{}_{ab}=0\ \text{ for all }a,b\in\{0,\dots,3\}
\]
(minus the $\Gamma$ symbol of lecture 7, \cref{prop:c07-fix}): ``time flows uniformly''
\ts{9}{1}{1:19:26}\ts{9}{1}{1:19:47}\ts{9}{1}{1:20:07}\ts{9}{1}{1:20:28}\ts{9}{1}{1:20:49}.

\section{Newton 2 in a stratified chart: the pseudo-accelerations}\label{sec:c09-newton2}

``Everything is said; what I am going to do now reconnects you with chart physics which we are used to.'' Newton 2 says $\nabla_VV=F/m$; take it in a chart $(U,x)$ of a stratified atlas, with the world line $X$ parametrised by some $\lambda$ (the world line is a future-directed curve; nothing was said about the parametrisation). The zeroth component of the auto-parallel equation is $X^{0\prime\prime}+\Gamma^0{}_{ab}X^{a\prime}X^{b\prime}=0$, hence ($\Gamma^0{}_{ab}=0$) $X^{0\prime\prime}=0$. Do we lose information about time? No: one evaluates the absolute time function at the point, whatever parameter was chosen. For the three spatial components, splitting the sum with the convention that Greek indices run over $1,2,3$ (``I forgot terms: both zero, one zero and one spatial'')
\[
  X^{\alpha\prime\prime}+\Gamma^\alpha{}_{\gamma\delta}X^{\gamma\prime}X^{\delta\prime}+2\,\Gamma^\alpha{}_{0\gamma}X^{0\prime}X^{\gamma\prime}+\Gamma^\alpha{}_{00}\bigl(X^{0\prime}\bigr)^2=\frac{F^\alpha}{m},
\]
where the factor $2$ appears because the mixed terms occur twice if $\nabla$ is \emph{torsion free}, $\Gamma^\alpha{}_{0\gamma}=\Gamma^\alpha{}_{\gamma0}$, and the right-hand side has $F^0=0$ in a stratified atlas (since $\mathrm dt(F)=0$): ``this is Newton's second law written out; I bet you spend a lot of time in your lectures on classical mechanics deriving this too''
\ts{9}{1}{1:21:13}\ts{9}{1}{1:21:36}\ts{9}{1}{1:21:57}\ts{9}{1}{1:22:19}\ts{9}{1}{1:22:39}\ts{9}{1}{1:22:59}\ts{9}{1}{1:23:20}\ts{9}{1}{1:23:41}\ts{9}{1}{1:24:01}\ts{9}{1}{1:24:22}\ts{9}{1}{1:25:25}\ts{9}{1}{1:25:46}.

\begin{coursenote}[Two axioms forgotten, then added]
While writing the terms Schuller realises that he forgot axioms: ``I made a postulate on the absolute time ($\mathrm dt\neq0$), I made a postulate on the compatibility of $\nabla$ and $t$ ($\nabla\mathrm dt=0$), and the third one I forgot is a postulate on $\nabla$ alone: \emph{$\nabla$ is torsion free}. One condition on the absolute time, one on the covariant derivative and one compatibility condition between the two.'' \cref{def:c09-nst} contains all three \ts{9}{1}{1:24:22}\ts{9}{1}{1:24:44}\ts{9}{1}{1:25:04}.
\end{coursenote}

\paragraph{The time component and the parameter.} In a stratified atlas $X^{0\prime\prime}=0$, so $X^0(\lambda)=a\lambda+b$ with constants $a,b$: the time of a point along the curve is $t\circ X(\lambda)=X^0(\lambda)$, \emph{proportional to the parameter}. ``The zeroth component is not garbage; before we had it as the zeroth equation and it was trivial; but that was only true in a stratified atlas'': the spacetime diagrams of the first part were drawn with space repeating itself, which looks like this only in a stratified atlas; ``what I did was right, but now I understand why: I already used the global time function to define the spaces''
\ts{9}{1}{1:26:06}\ts{9}{1}{1:26:31}\ts{9}{1}{1:26:52}\ts{9}{1}{1:27:14}\ts{9}{1}{1:27:34}\ts{9}{1}{1:27:55}.

\begin{coursenote}[Student question: why do we not just recover $\Gamma=0$?]
A student: in the first part of the lecture the $\Gamma$ with zero up were set to zero by demanding the form of Newton's law, and now it comes out again as a result: is this circular? Schuller: no. The first half starts with what one always writes down and works to an understanding that it looks better, but \emph{it was caught in preconceptions about charts and about what a Newtonian spacetime is}, formulated vaguely and non-geometrically, always in charts. Then ``I made a cut: forget everything from before, I start now with this abstract definition, I reformulated both laws in the new language, and I derive what I had before from this reformulation''. One cannot cross-link the two derivations; in particular, there is \emph{no assumption on $\nabla$ being flat} anywhere \ts{9}{1}{1:28:16}\ts{9}{1}{1:28:36}\ts{9}{1}{1:28:57}\ts{9}{1}{1:29:17}\ts{9}{1}{1:29:38}\ts{9}{1}{1:29:59}.
\end{coursenote}

\paragraph{Parametrising by the absolute time.} By convention one parametrises the world line by the absolute time: $d/d\lambda=a\,d/dt$ with constant $a$, so each prime becomes a dot with a factor $a$; the double derivative and the two velocity factors give a common factor $a^2$ on the left, which one divides out (it is just the unit of time one uses, seconds or years; on the right-hand side it survives in the force law, ``usually absorbed in the force''). One obtains the spatial equation of motion
\[
  \ddot X^\alpha+\Gamma^\alpha{}_{\gamma\delta}\dot X^\gamma\dot X^\delta+2\,\Gamma^\alpha{}_{0\gamma}\dot X^\gamma+\Gamma^\alpha{}_{00}=\frac{F^\alpha}{m}
\]
[with the unit-of-time factor $a^2$ absorbed into $F^\alpha/m$]
\ts{9}{1}{1:30:19}\ts{9}{1}{1:30:40}\ts{9}{1}{1:31:00}\ts{9}{1}{1:31:22}\ts{9}{1}{1:31:42}\ts{9}{1}{1:32:03}.
\begin{coursenote}[Student question on the factor $a^2$]
A student sees $a^2$ emerging on the left but not on the right. Schuller: ``I forgot the square'' on the right-hand side: it is there, and one divides through; the acceleration he writes with a small $a$ is just a symbol for acceleration, used because the index was also called $a$
\ts{9}{1}{1:33:46}\ts{9}{1}{1:34:06}\ts{9}{1}{1:34:27}.
\end{coursenote}

\paragraph{Only the sum is an acceleration.} What one would usually call the acceleration is $\ddot X^\alpha$, which is \emph{not} the acceleration: ``all of this, the $\alpha$ component, is the one and only acceleration; none of these terms by itself is an acceleration, if any of the terms is missing this no longer transforms as a vector''. The acceleration is a vector, hence a tensor: if it has vanishing components in one chart it has vanishing components in any chart (here a chart of \emph{spacetime}). If one twists the chart in spacetime one gets a rotating coordinate system in space, which is one fixed chart in spacetime: ``it describes the view of a dancer who carries the coordinate system with him'' \ts{9}{1}{1:32:03}\ts{9}{1}{1:32:23}\ts{9}{1}{1:32:45}\ts{9}{1}{1:33:05}\ts{9}{1}{1:33:26}\ts{9}{1}{1:33:46}.

\paragraph{Gravity, flatness and pseudo-accelerations.} In the presence of gravity the force $F^\alpha/m$ (the little $f^\alpha$) contributes to the acceleration through $\Gamma^\alpha{}_{00}$, a term that cannot be transformed away; in the right chart this is the only non-zero $\Gamma$ (``you can make all the others zero''). In the \emph{absence of gravity} one can formulate this as: \emph{Newtonian spacetime is flat}; there is a coordinate system with all $\Gamma$'s zero and one gets $\ddot x^\alpha=F^\alpha/m$ ``in a system without gravity and in a very special coordinate system'' (``one can reformulate this fully geometrically, I spare you this'').
Now view the same physics in a chart twisted in spacetime, which one would intuitively call a rotating chart (it could also stretch, ``do all crazy stuff''): even though in the previous chart all $\Gamma$'s were zero, all kinds of $\Gamma$'s pop up (a stratified atlas is not affected). Then the terms mean:
\begin{itemize}
  \item $\Gamma^\alpha{}_{\gamma\delta}\dot x^\gamma\dot x^\delta$, only spatial indices: a chart change in space, e.g.\ polar coordinates instead of Cartesian ones, the correction terms of polar coordinates ``you did this a thousand times in classical mechanics'';
  \item the term $2\,\Gamma^\alpha{}_{0\gamma}\dot x^\gamma$, for a rotating chart, is proportional to the spatial velocity of the particle (with the angular frequency $\Omega$ in it): the \textbf{Coriolis pseudo-acceleration}; the factor two is because the spacetime has no torsion and the term appears twice (``look at Wikipedia how you get the two'');
  \item $\Gamma^\alpha{}_{00}$, which in a stratified chart is independent of velocity but depends on the rotation, is the \textbf{centrifugal pseudo-acceleration}.
\end{itemize}
They are called pseudo-accelerations ``because you feel something is not quite right about them as an acceleration, of course there is not, because neither this is an acceleration nor that is an acceleration: only this plus this plus this is an acceleration''. In a rotating system $\ddot x^\alpha$ is no longer zero for a particle that was before moving uniformly free, and the other terms are no longer zero either, but \emph{the sum remains zero} because it is the acceleration of the particle, ``totally unimpressed by your fantasy about what coordinate system in spacetime you have''
\ts{9}{1}{1:34:27}\ts{9}{1}{1:34:48}\ts{9}{1}{1:35:08}\ts{9}{1}{1:35:29}\ts{9}{1}{1:35:49}\ts{9}{1}{1:36:10}\ts{9}{1}{1:36:31}\ts{9}{1}{1:36:51}\ts{9}{1}{1:37:12}\ts{9}{1}{1:37:33}\ts{9}{1}{1:37:53}\ts{9}{1}{1:38:14}\ts{9}{1}{1:38:34}\ts{9}{1}{1:38:55}\ts{9}{1}{1:39:15}\ts{9}{1}{1:39:36}\ts{9}{1}{1:39:57}.
\begin{aside}[A metric on space is an inertial-frame fiction]
To say what rotation is (rather than rotation plus dilation) one would need a metric on each spatial section, which has not been introduced yet ``for a good reason'': the idea that each absolute space carries a metric is \emph{an inertial-system fiction}; in a non-inertial system the geometry of space is more complicated than a three-dimensional metric geometry (a one-form $\Omega$ for the rotation is needed beside the metric $g$) \ts{9}{1}{1:37:12}\ts{9}{1}{1:37:33}\ts{9}{1}{1:46:51}\ts{9}{1}{1:47:12}.
\end{aside}

\section{Questions at the end}\label{sec:c09-questions}

\begin{coursenote}[Torsion and auto-parallels]
A student asks why the torsion is assumed zero, since this did not come out of the first half. Schuller: in the first part one assumed first that all the $\Gamma$'s are zero and then was forced to see that some are not; that came from ``prejudices to a certain level of understanding'', and then ``to get rid of all the dead bodies we left on the side of the road'' (a ``pedagogic omission of truth'') the second part comes top down and is the definite formulation. He first thought of starting with the last section but ``that is somewhat abstract if you have not understood the idea before why you would even be prompted to go there''. The student objects that auto-parallels are impressed by torsion if one has a metric compatible connection (next lecture) and adds torsion; Schuller: it is a notoriously confusing statement whether torsion does or does not affect the auto-parallel equation; the antisymmetric part falls out of the equation by inspection, but the question is whether the \emph{symmetric} part already has a torsion contribution by the construction mechanism; ``it depends on how you define the connection starting from given structures; here we just postulate it and everything works out like we like it'' \ts{9}{1}{1:40:18}\ts{9}{1}{1:40:38}\ts{9}{1}{1:40:59}\ts{9}{1}{1:41:19}\ts{9}{1}{1:41:41}\ts{9}{1}{1:42:02}\ts{9}{1}{1:42:23}\ts{9}{1}{1:42:43}.
\end{coursenote}
In Einstein gravity one requires the connection to be torsion free and \emph{compatible with the Lorentzian metric} of spacetime instead of with the time function, which ``philosophically replaces the time function'': the sections it makes are no longer slices but like the cone (everything inside time-like, outside space-like). The Lorentzian metric is not the generalisation of the spatial metric (``you can think of it that way, but it is particularly non-insightful''): \emph{it is the generalisation of the absolute time function} \ts{9}{1}{1:43:04}\ts{9}{1}{1:43:24}\ts{9}{1}{1:43:45}.

\begin{coursenote}[Student question: the rocket]
``Suppose I am in a rocket and turn on and off the engine: the force depends on time, so how do I use Newton 2?'' Schuller: $F$ is a spatial vector field, so $\mathrm dt(F)=0$, but it is a vector field on spacetime and can change as one proceeds in spacetime: it just has no component pointing out of the equal-time surface, and at another point it can be different; switching the rocket off and on again is captured \ts{9}{1}{1:43:45}\ts{9}{1}{1:44:05}\ts{9}{1}{1:44:26}.
\end{coursenote}
\begin{coursenote}[Student question: the Euler--Lagrange equations]
A student has an idea of deriving the Euler--Lagrange equations from this formulation. Not before one understands how to establish a metric on the spatial sections, since the Lagrangian invariably has the form of the kinetic energy $\tfrac12m\,g(\dot x,\dot x)$ with a spatial metric $g$, without a square root: the Newtonian Lagrangian formalism works only on the spatial section (``if you put the square root you have reparametrisation invariance and get straight curves but not straight uniform curves; if you take it away you fix the parameter automatically and get straight uniform curves: that is the reason for the kinetic energy being like this''). It does not work in spacetime as such, because only in \emph{inertial systems}, coordinate frames in which the $\Gamma$'s vanish (non-accelerated, non-rotating), can one establish a spatial metric everywhere. In a rotating system one needs also the angular-frequency one-form: ``so really having a metric of space is an inertial-system fiction'', and in an inertial system one can then change coordinates to depend also on time (``the formulas of Lagrangian mechanics stay true'') and pick up the terms, ``but you do not really understand them'' \ts{9}{1}{1:44:47}\ts{9}{1}{1:45:07}\ts{9}{1}{1:45:28}\ts{9}{1}{1:45:49}\ts{9}{1}{1:46:09}\ts{9}{1}{1:46:30}\ts{9}{1}{1:46:51}\ts{9}{1}{1:47:12}\ts{9}{1}{1:47:33}.
\end{coursenote}
Schuller mentions that his lectures in German on classical mechanics (a second-year undergraduate lecture, about 40 hours, on the web) follow this approach and finally come to Lagrangian mechanics: see \emph{Theoretische Mechanik} in the reading box \ts{9}{1}{1:47:53}\ts{9}{1}{1:48:14}.

% ==================================================================
\section*{Key concepts}
\begin{keyconcepts}
\textbf{Retrieval questions.}
\begin{enumerate}
  \item State Newton 1 and 2. Why must the first be read as a measurement prescription, and why is gravity as a force a problem for it?
  \item What is the weak equivalence principle and how was it tested in the lecture? What does it allow one to cancel?
  \item What was Laplace's question and why is the answer no? Why can $\ddot x=f(x)$ not be an auto-parallel equation on space?
  \item What is a world line, and why does moving to spacetime make straight \emph{and uniform} motion into straight motion?
  \item Which $\Gamma$'s turn Newton's gravitational equation into an auto-parallel equation in spacetime? Which components of curvature are non-zero, and how does Poisson's equation appear?
  \item Why is the curvature not a coordinate artifact, and why is it spacetime and not space curvature?
  \item Define a Newtonian spacetime. How does one obtain absolute space, and what does $\nabla\mathrm dt=0$ say in a stratified atlas?
  \item What are the Coriolis and centrifugal terms in this language, and why is only their sum with $\ddot x$ an acceleration?
\end{enumerate}
\textbf{Concepts.}
\begin{itemize}
  \item Newton 1 as a measurement prescription; weak equivalence principle; Laplace's question (\cref{sec:c09-laplace}).
  \item World lines and $\Gamma^\alpha{}_{00}=-f^\alpha$ (\cref{prop:c09-gamma}); tidal tensor $R^\alpha{}_{0\beta0}=-\partial_\beta f^\alpha$; $R_{00}=4\pi G\rho$.
  \item Newtonian spacetime $(M,\mathcal O,\mathcal A,\nabla,t)$ (\cref{def:c09-nst}); future-directed vectors (\cref{def:c09-fut}); Newton 1 and 2 geometrically (\cref{def:c09-newtongeom}).
  \item Stratified atlas; $\Gamma^0{}_{ab}=0$; the acceleration is the sum $\ddot x^\alpha+\Gamma\dot x\dot x+2\Gamma\dot x+\Gamma$.
\end{itemize}
\end{keyconcepts}

\flushside
\end{paracol}
